\begin{solapartat}
\punts{0,2} Equilibri de forces sobre la partícula, $\Sigma F = 0$

\smallskip
\punts{0,4} $F_e = F_g \Rightarrow qE = mg\,;\ \dfrac{q}{m} = \dfrac{g}{E}$

\smallskip
\punts{0,2} $\dfrac{q}{m} = \dfrac{1,62}{10} = 0,162\ \mathrm{C\,kg^{-1}}$

\smallskip
\punts{0,2} La càrrega de la partícula \textbf{és positiva.}
\end{solapartat}

\begin{solapartat}
A la cara fosca de la Lluna, segons l'enunciat, la força gravitatòria i la força elèctrica sobre una càrrega positiva tenen el mateix sentit, cap al centre de la Lluna. A partir d'aquí, considerarem amb signe positiu els vectors que es dirigeixen cap a l'exterior. Únicament hi intervenen les coordenades verticals.

\medskip
\punts{0,2} $\Sigma F = mg + qE$

\smallskip
\punts{0,1} $\Sigma F = 0,020\times 10^{-6}(1,62) + 10\times 10^{-9}(1,0) = 4,2\times 10^{-8}\ \mathrm{N}$ (mòdul)

\smallskip
\punts{0,1} Aquesta força resultant va \textbf{dirigida} cap al centre de la Lluna, segons la direcció i sentit del camp elèctric.

En 10 metres no canviaran de manera apreciable $\vec{g}$ i $\vec{E}$.

\smallskip
\punts{0,2} $a = \dfrac{\Sigma F}{m} = \dfrac{4,24\times 10^{-8}}{0,020\times 10^{-6}} = 2,12\ \mathrm{m\,s^{-2}}$ cap al centre de la Lluna

\smallskip
\punts{0,2} $y - y_0 = v_{y0}t + \dfrac{1}{2}a_y t^2$

\smallskip
\punts{0,2} $t = \sqrt{\dfrac{10}{1,06}} = 3,07\ \mathrm{s}$
\end{solapartat}
